\documentclass[10pt,a4paper]{article}
\usepackage[utf8]{inputenc}
\usepackage{amsmath}
\usepackage{amsfonts}
\usepackage{amssymb}
\usepackage{enumerate}
\title{GSF-3100 \\
03 Volatilité dans le marché obligataire \\
Série d'exercices 4}

\begin{document}
\maketitle

\begin{itemize}
\item Dans ce document vous avez des questions à choix multiples allant de a) à d).  
\item Il y a seulement un choix possible par question.
\item Ce questionnaire contient des exercices en lien avec la volatilité dans le marché obligataire.
\item Contient 9 questions. 
\end{itemize}
 





\newpage

\section*{Exercice 1}
Pour une obligation à valeur nominale M = 100 \$ venant à échéance dans 60 semestres avec un taux de coupon annuel $TC= 3$ \% payable quatre fois par année et un taux de rendement effectif trimestriel exigé de y = 1 \%,  trouver la durée modifiée en année.


\begin{enumerate}[a)]
\item 18,50142 
\item 18,60142 
\item 18,55142 
\item 18,65142 
\end{enumerate}

\section*{Exercice 2}
Pour une obligation à valeur nominale M = 100 \$ venant à échéance dans 60 semestres avec un taux de coupon annuel $TC= 3$ \% payable quatre fois par année et un taux de rendement effectif trimestriel exigé de y = 1 \%,  trouver la durée de Macaulay en trimestre.


\begin{enumerate}[a)]
\item 74.94773 
\item 74.99773 
\item 74.90773 
\item 74.94073 
\end{enumerate}

\section*{Exercice 3}
Considérer l’obligation suivante: TC effectif semestriel = 2 \%, n = 20 semestres, y effectif semestriel = 1,6 \%, P = 106,80023 \$, M = 100 \$ et $D_m$ en année = 7,490295 années.  Si on applique une diminution du taux de rendement de 150 points de base,  soit $\Delta y = -1,5$ \%,  trouver la vraie variation de prix $P_{y+\Delta y} - P$ en dollar.


\begin{enumerate}[a)]
\item 23.99896 \$ 
\item 30.80368 \$
\item 31.80368 \$
\item  24.99896 \$
\end{enumerate}

\newpage

\section*{Exercice 4}
Considérer l’obligation suivante: TC effectif semestriel = 2 \%, n = 20 semestres, y effectif semestriel = 1,6 \%, P = 106,80023 \$, M = 100 \$ et $D_m$ en année = 7,490295 années.  Si on applique une diminution du taux de rendement de 150 points de base,  soit $\Delta y = -1,5$ \%,  trouver ;a variation approximative du prix $\Delta P$ en dollar.


\begin{enumerate}[a)]
\item 23.99896 \$
\item 30.80368 \$
\item 31.80368 \$
\item  24.99896 \$
\end{enumerate}


\section*{Exercice 5}
Pour une obligation à valeur nominale $M = 100$ \$ venant à échéance dans $n = 60$ semestres avec un taux de coupon annuel $TC = 0$ \% payable quatre fois par année et un taux de rendement effectif trimestriel exigé de $y = 1$ \%,  trouver la durée de Macaulay en année.


\begin{enumerate}[a)]
\item 29.70297 
\item 30.0000
\item 30.5000
\item 30.70297 
\end{enumerate}


\section*{Exercice 6}
Pour une obligation à valeur nominale $M = 100$ \$ venant à échéance dans $n = 60$ semestres avec un taux de coupon annuel $TC = 0$ \% payable quatre fois par année et un taux de rendement effectif trimestriel exigé de $y = 1$ \%,  trouver la durée modifiée en année.

\begin{enumerate}[a)]
\item 30.5000
\item 30.0000
\item 29.70297 
\item 30.70297 
\end{enumerate}

\newpage

\section*{Exercice 7}
Considérer l’obligation suivante: TC effectif périodique = 0 \%, $n = 60$, y effectif périodique = 0,5 \%, $P = 7413,72196$ \$, $M= 10000$ \$, $D_m = 59.70149$ et $C = 3623.67268$.  Si on augmente le taux de rendement de 200 points de base ($\Delta y = 2$  \%),  trouver la variation de prix en pourcentage en utilisant la durée et la convexité (approximation quadratique). 

\begin{enumerate}[a)]
\item -119.403 \%
\item -46.930 \%
\item -26.930 \%
\item -69.343 \%
\end{enumerate}

\section*{Exercice 8}
Considérer l’obligation suivante: TC effectif périodique = 0 \%, $n = 60$, y effectif périodique = 0,5 \%, $P = 7413,72196$ \$, $M= 10000$ \$, $D_m = 59.70149$ et $C = 3623.67268$.  Si on augmente le taux de rendement de 200 points de base ($\Delta y = 2$  \%),  trouver la variation de prix en pourcentage en utilisant la durée seulement (approximation linéaire).

\begin{enumerate}[a)]
\item -119.403 \%
\item -46.930 \%
\item -26.930 \%
\item -69.343 \%
\end{enumerate}


\section*{Exercice 9}
Considérer l’obligation suivante: TC effectif périodique = 0 \%, $n = 60$, y effectif périodique = 0,5 \%, $P = 7413,72196$ \$, $M= 10000$ \$, $D_m = 59.70149$ et $C = 3623.67268$.  Si on augmente le taux de rendement de 200 points de base ($\Delta y = 2$  \%),  trouver la variation de prix en pourcentage en utilisant le vrai prix après le changement de taux $P_{y+\Delta y}$.

\begin{enumerate}[a)]
\item -119.403 \%
\item -46.930 \%
\item -26.930 \%
\item -69.343 \%
\end{enumerate}


\end{document}